\section{Results}
\label{sec:results}

\subsection{Memory frontier: where dense training stops}
\label{sec:memory}

% Source: docs/GRPO_CONTEXT_SCALING.md (Qwen2.5-7B-Instruct, L40S 48 GB,
% group 2, 32 new tokens, chunk 1024, 1 iter). Regenerate with
% examples/grpo_context_sweep.py.
\begin{table}[t]
  \caption{Memory frontier for a single GRPO step, Qwen2.5-7B-Instruct on one
  48\,GB L40S. Dense attention runs out of memory at 24k context and a
  full-coverage sparse cache at 32k; a bounded evicting cache (budget 8192)
  is flat at $\sim$35.7\,GB up to 65k. Peak is
  \texttt{torch.cuda.max\_memory\_allocated}; times are CUDA-synchronised
  wall clock for generation and the chunked gradient pass.}
  \label{tab:memory_frontier}
  \centering\scriptsize
  \setlength{\tabcolsep}{3pt}
  \begin{tabular}{rlrrrr}
    \toprule
    Context & Cache & Gen (s) & Grad (s) & Total (s) & Peak (GB) \\
    \midrule
    8{,}192  & dense              & 4.4  & 8.7   & 13.4  & 35.0 \\
    16{,}384 & dense              & 7.8  & 18.1  & 26.0  & 40.8 \\
    24{,}576 & dense              & --   & --    & OOM   & OOM \\
    \midrule
    16{,}384 & H2O-q8, full & 9.1  & 18.4  & 28.0  & 40.9 \\
    32{,}768 & H2O-q8, full & --   & --    & OOM   & OOM \\
    \midrule
    24{,}576 & H2O-q8, $K{=}$8k  & 20.1 & 54.5  & 75.1  & 35.5 \\
    32{,}768 & H2O-q8, $K{=}$8k  & 27.5 & 76.9  & 104.7 & 35.7 \\
    49{,}152 & H2O-q8, $K{=}$8k  & 41.8 & 120.7 & 162.9 & 35.7 \\
    65{,}536 & H2O-q8, $K{=}$8k  & 56.2 & 163.8 & 220.9 & 35.7 \\
    \midrule
    24{,}576 & H2O-q8, $K{=}$12k & 22.7 & 64.6  & 87.9  & 38.5 \\
    32{,}768 & H2O-q8, $K{=}$12k & 32.3 & 95.2  & 127.7 & 38.5 \\
    \bottomrule
  \end{tabular}
\end{table}


\cref{tab:memory_frontier} gives the single-step footprint for
Qwen2.5-7B-Instruct on a 48\,GB L40S. Dense GRPO fails at 24k. A sparse cache
large enough to cover the whole prompt fails at 32k because the cache is paid
for again in the gradient pass. A bounded evicting cache with $K=8192$ holds
peak memory at 35.7\,GB from 24k to 65k context, $2.7\times$ past the dense
limit. The price is step time: the chunked gradient pass is linear in prompt
length and the cache adds a scoring and eviction cost, so a step at 32k takes
$\sim$105\,s and at 65k $\sim$221\,s (group of 2). Where both fit (${\le}16$k)
the evicting cache is $1.3$--$1.5\times$ slower than dense. \draft{Add the
3B/A10G table to the appendix as a second data point (dense OOM at 24k on
24\,GB, H2O flat at 14.9\,GB to 32k).}

\subsection{Decode throughput: where dense generation slows down}
\label{sec:decode}

% Source: docs/DECODE_BENCH.md (run aa0_decode_bench_7b, 2026-10-01).
% Qwen2.5-7B-Instruct, L40S 48 GB, batch 8, 256 decoded tokens, eager SDPA.
\begin{table}[t]
  \caption{Decode throughput (tokens/s, batch 8) and peak memory for
  Qwen2.5-7B-Instruct on one 48\,GB L40S. Dense decode slows as the context
  grows and runs out of memory at 16k; a bounded H2O cache ($K{=}8192$) is
  flat in both throughput and memory from 8k to 65k. The int8 KV cache
  trades a third of the throughput for 1.5\,GB.}
  \label{tab:decode_throughput}
  \centering\scriptsize
  \setlength{\tabcolsep}{3pt}
  \begin{tabular}{rrrrrrr}
    \toprule
    & \multicolumn{2}{c}{dense} & \multicolumn{2}{c}{H2O bf16, $K{=}$8k} & \multicolumn{2}{c}{H2O int8, $K{=}$8k} \\
    \cmidrule(lr){2-3}\cmidrule(lr){4-5}\cmidrule(lr){6-7}
    Context & tok/s & GB & tok/s & GB & tok/s & GB \\
    \midrule
    4{,}096  & 59.1 & 21.8 & 58.3 & 21.8 & 47.6 & 21.0 \\
    8{,}192  & 32.4 & 29.3 & 32.3 & 29.3 & 22.5 & 27.7 \\
    16{,}384 & OOM  & OOM  & 32.3 & 29.3 & 21.9 & 27.7 \\
    32{,}768 & OOM  & OOM  & 32.1 & 29.3 & 21.9 & 27.8 \\
    \bottomrule
  \end{tabular}
\end{table}


Rollout generation is memory-bandwidth bound: every decoded token reads the
full KV cache of every layer, so per-token cost grows with the context.
\cref{tab:decode_throughput} measures this directly for a batch of 8
sequences (one GRPO group) decoding 256 tokens. Dense throughput halves from
4k to 8k and the batch no longer fits at 16k. The bounded cache reads $K$
slots per token regardless of context: from 8k to 32k its throughput and
memory do not move. Where both fit, the bf16 evicting cache matches dense
within 1\%, so the eviction bookkeeping (score accumulation and slot
ranking) is not a measurable cost at this scale; amortising the ranking
over blocks of 64 tokens changes nothing. The int8 KV cache, which is the
right choice for the gradient pass (\cref{sec:gap}), costs a third of the
decode throughput in dequantisation and write-back, so we use bf16 for
rollouts and int8 where memory, not time, is binding. \draft{Add the
profiler breakdown as a sentence or footnote: 27\% of CUDA time in copies
for int8; the remaining fixed overhead is the eager attention path that
returns per-slot weights for H2O scoring.}

\subsection{Parity where dense fits}
\label{sec:parity}

\draft{Shortest generation cutoff at which dense GRPO still fits on one
L40S for Qwen3-1.7B (and 4B if possible). GRPO and RLOO, dense vs evicting
cache, 2 seeds. Table: AIME24/25/26, AMC23/24, MATH500 accuracy plus training
reward at the end of one epoch on Polaris. Headline we want: every paired
comparison within noise at lower memory. The HELMET 0.5B@8k campaign
(docs/GRPO\_HELMET\_RESULTS.md) is the existing evidence for this claim on a
different task family; move it to the appendix as supporting data.}

\subsection{Gains where dense cannot run}
\label{sec:flagship}

\draft{Generation cutoff past the dense OOM point (32k on 48\,GB for the
larger model). Only the bounded evicting cache trains here. Report
through-cache and dense-eval accuracy on the math suite before and after one
epoch on Polaris, and the training reward curve next to Sparrow's
Figure~6(a) convention. Question to answer explicitly: does training through
the cache hurt the dense-eval numbers?}

\subsection{Coding RL}
\label{sec:coding}

\draft{Qwen3-1.7B on 21K filtered TACO, sandboxed unit-test reward, same
cache config as the math runs. Evaluate LiveCodeBench, HumanEval+, MBPP+.
Mirrors Sparrow Table~3 so the reader can line the two up.}

\subsection{What eviction destroys and what RL recovers}
\label{sec:gap}

\draft{Fold in the A1/A2 analyses (docs \texttt{POSITION\_COMPACTION\_A1} and
\texttt{EVICTION\_SWEEP\_A2}). The gap is information loss:
(i) exact RoPE re-indexing (position compaction) recovers nothing;
(ii) no retention knob at fixed slots recovers anything;
(iii) 8-bit KV quantisation is accuracy-free (dense-q8 $=$ h2o-q8 controls),
so at equal memory you afford $2\times$ slots and, when the cache covers the
prompt, accuracy equals dense at half the memory;
(iv) under genuine eviction the gap degrades gracefully, $\sim$25\% retention
costs $\sim$9\,pp EM on trivia QA (HELMET, appendix). Then: RL through the cache can only teach the
policy to use what survives; quantify how much of the gap is recoverable.}

\subsection{Ablations}
\label{sec:ablations}

\draft{Candidates: cache budget $K$ (8192 vs 12288: $+3$\,GB, $+20\%$ step
time, 50\% more prompt retained); GRPO vs RLOO; DeepScaleR warm-up on/off; shared-prompt prefill on/off (speed only, token-exact).}
